%!TeX root=../princesstop.tex
\chapter{Melchisedec}

\begin{figure}[t!]
\centering
\includegraphics[width=\linewidth]{filigree}
\end{figure}

\lettrine[lines=5]{T}{he} third person in the trio was Lottie. She was a small thing and did not know what adversity meant, and was much bewildered by the alteration she saw in her young adopted mother. She had heard it rumoured that strange things had happened to Sara, but she could not understand why she looked different—why she wore an old black frock and came into the schoolroom only to teach instead of to sit in her place of honour and learn lessons herself. There had been much whispering among the little ones when it had been discovered that Sara no longer lived in the rooms in which Emily had so long sat in state. Lottie's chief difficulty was that Sara said so little when one asked her questions. At seven mysteries must be made very clear if one is to understand them.

<Are you very poor now, Sara?> she had asked confidentially the first morning her friend took charge of the small French class. <Are you as poor as a beggar?> She thrust a fat hand into the slim one and opened round, tearful eyes. <I don't want you to be as poor as a beggar.>

She looked as if she was going to cry. And Sara hurriedly consoled her.

<Beggars have nowhere to live,> she said courageously. <I have a place to live in.>

<Where do you live?> persisted Lottie. <The new girl sleeps in your room, and it isn't pretty any more.>

<I live in another room,> said Sara.

<Is it a nice one?> inquired Lottie. <I want to go and see it.>

<You must not talk,> said Sara. <Miss Minchin is looking at us. She will be angry with me for letting you whisper.>

She had found out already that she was to be held accountable for everything which was objected to. If the children were not attentive, if they talked, if they were restless, it was she who would be reproved.

But Lottie was a determined little person. If Sara would not tell her where she lived, she would find out in some other way. She talked to her small companions and hung about the elder girls and listened when they were gossiping; and acting upon certain information they had unconsciously let drop, she started late one afternoon on a voyage of discovery, climbing stairs she had never known the existence of, until she reached the attic floor. There she found two doors near each other, and opening one, she saw her beloved Sara standing upon an old table and looking out of a window.

<Sara!> she cried, aghast. <Mamma Sara!> She was aghast because the attic was so bare and ugly and seemed so far away from all the world. Her short legs had seemed to have been mounting hundreds of stairs.

Sara turned round at the sound of her voice. It was her turn to be aghast. What would happen now? If Lottie began to cry and any one chanced to hear, they were both lost. She jumped down from her table and ran to the child.

<Don't cry and make a noise,> she implored. <I shall be scolded if you do, and I have been scolded all day. It's—it's not such a bad room, Lottie.>

<Isn't it?> gasped Lottie, and as she looked round it she bit her lip. She was a spoiled child yet, but she was fond enough of her adopted parent to make an effort to control herself for her sake. Then, somehow, it was quite possible that any place in which Sara lived might turn out to be nice. <Why isn't it, Sara?> she almost whispered.

Sara hugged her close and tried to laugh. There was a sort of comfort in the warmth of the plump, childish body. She had had a hard day and had been staring out of the windows with hot eyes.

<You can see all sorts of things you can't see downstairs,> she said.

<What sort of things?> demanded Lottie, with that curiosity Sara could always awaken even in bigger girls.

<Chimneys—quite close to us—with smoke curling up in wreaths and clouds and going up into the sky—and sparrows hopping about and talking to each other just as if they were people—and other attic windows where heads may pop out any minute and you can wonder who they belong to. And it all feels as high up—as if it was another world.>

<Oh, let me see it!> cried Lottie. <Lift me up!>

Sara lifted her up, and they stood on the old table together and leaned on the edge of the flat window in the roof, and looked out.

Anyone who has not done this does not know what a different world they saw. The slates spread out on either side of them and slanted down into the rain gutter-pipes. The sparrows, being at home there, twittered and hopped about quite without fear. Two of them perched on the chimney top nearest and quarrelled with each other fiercely until one pecked the other and drove him away. The garret window next to theirs was shut because the house next door was empty.

<I wish someone lived there,> Sara said. <It is so close that if there was a little girl in the attic, we could talk to each other through the windows and climb over to see each other, if we were not afraid of falling.>

The sky seemed so much nearer than when one saw it from the street, that Lottie was enchanted. From the attic window, among the chimney pots, the things which were happening in the world below seemed almost unreal. One scarcely believed in the existence of Miss Minchin and Miss Amelia and the schoolroom, and the roll of wheels in the square seemed a sound belonging to another existence.

<Oh, Sara!> cried Lottie, cuddling in her guarding arm. <I like this attic—I like it! It is nicer than downstairs!>

<Look at that sparrow,> whispered Sara. <I wish I had some crumbs to throw to him.>

<I have some!> came in a little shriek from Lottie. <I have part of a bun in my pocket; I bought it with my penny yesterday, and I saved a bit.>

When they threw out a few crumbs the sparrow jumped and flew away to an adjacent chimney top. He was evidently not accustomed to intimates in attics, and unexpected crumbs startled him. But when Lottie remained quite still and Sara chirped very softly—almost as if she were a sparrow herself—he saw that the thing which had alarmed him represented hospitality, after all. He put his head on one side, and from his perch on the chimney looked down at the crumbs with twinkling eyes. Lottie could scarcely keep still.

<Will he come? Will he come?> she whispered.

<His eyes look as if he would,> Sara whispered back. <He is thinking and thinking whether he dare. Yes, he will! Yes, he is coming!>

He flew down and hopped toward the crumbs, but stopped a few inches away from them, putting his head on one side again, as if reflecting on the chances that Sara and Lottie might turn out to be big cats and jump on him. At last his heart told him they were really nicer than they looked, and he hopped nearer and nearer, darted at the biggest crumb with a lightning peck, seized it, and carried it away to the other side of his chimney.

<Now he \textsc{knows}>, said Sara. <And he will come back for the others.>

He did come back, and even brought a friend, and the friend went away and brought a relative, and among them they made a hearty meal over which they twittered and chattered and exclaimed, stopping every now and then to put their heads on one side and examine Lottie and Sara. Lottie was so delighted that she quite forgot her first shocked impression of the attic. In fact, when she was lifted down from the table and returned to earthly things, as it were, Sara was able to point out to her many beauties in the room which she herself would not have suspected the existence of.

<It is so little and so high above everything,> she said, <that it is almost like a nest in a tree. The slanting ceiling is so funny. See, you can scarcely stand up at this end of the room; and when the morning begins to come I can lie in bed and look right up into the sky through that flat window in the roof. It is like a square patch of light. If the sun is going to shine, little pink clouds float about, and I feel as if I could touch them. And if it rains, the drops patter and patter as if they were saying something nice. Then if there are stars, you can lie and try to count how many go into the patch. It takes such a lot. And just look at that tiny, rusty grate in the corner. If it was polished and there was a fire in it, just think how nice it would be. You see, it's really a beautiful little room.>

She was walking round the small place, holding Lottie's hand and making gestures which described all the beauties she was making herself see. She quite made Lottie see them, too. Lottie could always believe in the things Sara made pictures of.

<You see,> she said, <there could be a thick, soft blue Indian rug on the floor; and in that corner there could be a soft little sofa, with cushions to curl up on; and just over it could be a shelf full of books so that one could reach them easily; and there could be a fur rug before the fire, and hangings on the wall to cover up the whitewash, and pictures. They would have to be little ones, but they could be beautiful; and there could be a lamp with a deep rose-coloured shade; and a table in the middle, with things to have tea with; and a little fat copper kettle singing on the hob; and the bed could be quite different. It could be made soft and covered with a lovely silk coverlet. It could be beautiful. And perhaps we could coax the sparrows until we made such friends with them that they would come and peck at the window and ask to be let in.>

<Oh, Sara!> cried Lottie. <I should like to live here!>

When Sara had persuaded her to go downstairs again, and, after setting her on her way, had come back to her attic, she stood in the middle of it and looked about her. The enchantment of her imaginings for Lottie had died away. The bed was hard and covered with its dingy quilt. The whitewashed wall showed its broken patches, the floor was cold and bare, the grate was broken and rusty, and the battered footstool, tilted sideways on its injured leg, the only seat in the room. She sat down on it for a few minutes and let her head drop in her hands. The mere fact that Lottie had come and gone away again made things seem a little worse—just as perhaps prisoners feel a little more desolate after visitors come and go, leaving them behind.

<It's a lonely place,> she said. <Sometimes it's the loneliest place in the world.>

She was sitting in this way when her attention was attracted by a slight sound near her. She lifted her head to see where it came from, and if she had been a nervous child she would have left her seat on the battered footstool in a great hurry. A large rat was sitting up on his hind quarters and sniffing the air in an interested manner. Some of Lottie's crumbs had dropped upon the floor and their scent had drawn him out of his hole.

He looked so queer and so like a gray-whiskered dwarf or gnome that Sara was rather fascinated. He looked at her with his bright eyes, as if he were asking a question. He was evidently so doubtful that one of the child's queer thoughts came into her mind.

<I dare say it is rather hard to be a rat,> she mused. <Nobody likes you. People jump and run away and scream out, <Oh, a horrid rat!> I shouldn't like people to scream and jump and say, <Oh, a horrid Sara!> the moment they saw me. And set traps for me, and pretend they were dinner. It's so different to be a sparrow. But nobody asked this rat if he wanted to be a rat when he was made. Nobody said, <Wouldn't you rather be a sparrow?>>

She had sat so quietly that the rat had begun to take courage. He was very much afraid of her, but perhaps he had a heart like the sparrow and it told him that she was not a thing which pounced. He was very hungry. He had a wife and a large family in the wall, and they had had frightfully bad luck for several days. He had left the children crying bitterly, and felt he would risk a good deal for a few crumbs, so he cautiously dropped upon his feet.

<Come on,> said Sara; <I'm not a trap. You can have them, poor thing! Prisoners in the Bastille used to make friends with rats. Suppose I make friends with you.>

How it is that animals understand things I do not know, but it is certain that they do understand. Perhaps there is a language which is not made of words and everything in the world understands it. Perhaps there is a soul hidden in everything and it can always speak, without even making a sound, to another soul. But whatsoever was the reason, the rat knew from that moment that he was safe—even though he was a rat. He knew that this young human being sitting on the red footstool would not jump up and terrify him with wild, sharp noises or throw heavy objects at him which, if they did not fall and crush him, would send him limping in his scurry back to his hole. He was really a very nice rat, and did not mean the least harm. When he had stood on his hind legs and sniffed the air, with his bright eyes fixed on Sara, he had hoped that she would understand this, and would not begin by hating him as an enemy. When the mysterious thing which speaks without saying any words told him that she would not, he went softly toward the crumbs and began to eat them. As he did it he glanced every now and then at Sara, just as the sparrows had done, and his expression was so very apologetic that it touched her heart.

She sat and watched him without making any movement. One crumb was very much larger than the others—in fact, it could scarcely be called a crumb. It was evident that he wanted that piece very much, but it lay quite near the footstool and he was still rather timid.

<I believe he wants it to carry to his family in the wall,> Sara thought. <If I do not stir at all, perhaps he will come and get it.>

She scarcely allowed herself to breathe, she was so deeply interested. The rat shuffled a little nearer and ate a few more crumbs, then he stopped and sniffed delicately, giving a side glance at the occupant of the footstool; then he darted at the piece of bun with something very like the sudden boldness of the sparrow, and the instant he had possession of it fled back to the wall, slipped down a crack in the skirting board, and was gone.

<I knew he wanted it for his children,> said Sara. <I do believe I could make friends with him.>

A week or so afterward, on one of the rare nights when Ermengarde found it safe to steal up to the attic, when she tapped on the door with the tips of her fingers Sara did not come to her for two or three minutes. There was, indeed, such a silence in the room at first that Ermengarde wondered if she could have fallen asleep. Then, to her surprise, she heard her utter a little, low laugh and speak coaxingly to someone.

<There!> Ermengarde heard her say. <Take it and go home, Melchisedec! Go home to your wife!>

Almost immediately Sara opened the door, and when she did so she found Ermengarde standing with alarmed eyes upon the threshold.

<Who—who \textsc{are} you talking to, Sara?> she gasped out.

Sara drew her in cautiously, but she looked as if something pleased and amused her.

<You must promise not to be frightened—not to scream the least bit, or I can't tell you,> she answered.

Ermengarde felt almost inclined to scream on the spot, but managed to control herself. She looked all round the attic and saw no one. And yet Sara had certainly been speaking \textsc{to} someone. She thought of ghosts.

<Is it—something that will frighten me?> she asked timorously.

<Some people are afraid of them,> said Sara. <I was at first—but I am not now.>

<Was it—a ghost?> quaked Ermengarde.

<No,> said Sara, laughing. <It was my rat.>

Ermengarde made one bound, and landed in the middle of the little dingy bed. She tucked her feet under her nightgown and the red shawl. She did not scream, but she gasped with fright.

<Oh! Oh!> she cried under her breath. <A rat! A rat!>

<I was afraid you would be frightened,> said Sara. <But you needn't be. I am making him tame. He actually knows me and comes out when I call him. Are you too frightened to want to see him?>

The truth was that, as the days had gone on and, with the aid of scraps brought up from the kitchen, her curious friendship had developed, she had gradually forgotten that the timid creature she was becoming familiar with was a mere rat.

At first Ermengarde was too much alarmed to do anything but huddle in a heap upon the bed and tuck up her feet, but the sight of Sara's composed little countenance and the story of Melchisedec's first appearance began at last to rouse her curiosity, and she leaned forward over the edge of the bed and watched Sara go and kneel down by the hole in the skirting board.

<He—he won't run out quickly and jump on the bed, will he?> she said.

<No,> answered Sara. <He's as polite as we are. He is just like a person. Now watch!>

She began to make a low, whistling sound—so low and coaxing that it could only have been heard in entire stillness. She did it several times, looking entirely absorbed in it. Ermengarde thought she looked as if she were working a spell. And at last, evidently in response to it, a gray-whiskered, bright-eyed head peeped out of the hole. Sara had some crumbs in her hand. She dropped them, and Melchisedec came quietly forth and ate them. A piece of larger size than the rest he took and carried in the most businesslike manner back to his home.

<You see,> said Sara, <that is for his wife and children. He is very nice. He only eats the little bits. After he goes back I can always hear his family squeaking for joy. There are three kinds of squeaks. One kind is the children's, and one is Mrs Melchisedec's, and one is Melchisedec's own.>

Ermengarde began to laugh.

<Oh, Sara!> she said. <You \textsc{are} queer—but you are nice.>

<I know I am queer,> admitted Sara, cheerfully; <and I \textsc{try} to be nice.> She rubbed her forehead with her little brown paw, and a puzzled, tender look came into her face. <Papa always laughed at me,> she said; <but I liked it. He thought I was queer, but he liked me to make up things. I—I can't help making up things. If I didn't, I don't believe I could live.> She paused and glanced around the attic. <I'm sure I couldn't live here,> she added in a low voice.

Ermengarde was interested, as she always was. <When you talk about things,> she said, <they seem as if they grew real. You talk about Melchisedec as if he was a person.>

<He \textsc{is} a person,> said Sara. <He gets hungry and frightened, just as we do; and he is married and has children. How do we know he doesn't think things, just as we do? His eyes look as if he was a person. That was why I gave him a name.>

She sat down on the floor in her favourite attitude, holding her knees.

<Besides,> she said, <he is a Bastille rat sent to be my friend. I can always get a bit of bread the cook has thrown away, and it is quite enough to support him.>

<Is it the Bastille yet?> asked Ermengarde, eagerly. <Do you always pretend it is the Bastille?>

<Nearly always,> answered Sara. <Sometimes I try to pretend it is another kind of place; but the Bastille is generally easiest—particularly when it is cold.>

Just at that moment Ermengarde almost jumped off the bed, she was so startled by a sound she heard. It was like two distinct knocks on the wall.

<What is that?> she exclaimed.

Sara got up from the floor and answered quite dramatically:

<It is the prisoner in the next cell.>

<Becky!> cried Ermengarde, enraptured.

<Yes,> said Sara. <Listen; the two knocks meant, <Prisoner, are you there?>>

She knocked three times on the wall herself, as if in answer.

<That means, <Yes, I am here, and all is well.>>

Four knocks came from Becky's side of the wall.

<That means,> explained Sara, <<Then, fellow-sufferer, we will sleep in peace. Good night.>>

Ermengarde quite beamed with delight.

<Oh, Sara!> she whispered joyfully. <It is like a story!>

<It \textsc{is} a story,> said Sara. <\textsc{everything's} a story. You are a story—I am a story. Miss Minchin is a story.>

And she sat down again and talked until Ermengarde forgot that she was a sort of escaped prisoner herself, and had to be reminded by Sara that she could not remain in the Bastille all night, but must steal noiselessly downstairs again and creep back into her deserted bed.